\section{Related Work}
\label{sec:related}

\subsection{Robotic Garment Manipulation}
Manipulating highly deformable textiles presents unique challenges compared to rigid objects due to infinite-dimensional configuration spaces and complex dynamics \cite{longhini2024unfolding, 8957044}. Recent frameworks have automated tasks like smoothing and folding by employing dense visual correspondences, learned keypoints, and deep imitation learning \cite{ganapathi2020learningdensevisualcorrespondences, grannen2020untanglingdenseknotslearning}. Systems like SpeedFolding \cite{avigal2022speedfolding} and UniFolding \cite{unifolding2023} leverage bimanual coordination for efficient folding, often using discrete primitives or canonical alignments. Our work explores the potential of continuous action spaces modeled by generative architectures.

\subsection{Visuomotor Imitation Learning}
Imitation learning has shifted from traditional behavioral cloning to generative models that handle multi-modality. Diffusion Policy \cite{chi2023diffusionpolicy} models the action distribution via a diffusion process, while Action Chunking with Transformers (ACT) \cite{zhao2023act} leverages the temporal modeling of Transformers to predict chunks of actions, mitigating compounding errors. Recent advancements have introduced Large Vision-Language-Action (VLA) models such as RT-2 \cite{brohan2023rt2visionlanguageactionmodelstransfer}, OpenVLA \cite{kim2024openvlaopensourcevisionlanguageactionmodel}, and the $\pi$ series ($\pi_0$, $\pi_{0.5}$) \cite{black2026pi0visionlanguageactionflowmodel, pi05_2026}. These models often use text tokens to encode actions and leverage web-scale knowledge for robotic control. We investigate how these foundation models adapt to the specialized domain of garment folding.

\subsection{Generalization and Adaptation}
Achieving robust generalization across object styles requires effective visual representations and efficient adaptation techniques. 3D Diffusion Policy (DP3) \cite{ze2024dp3} uses point clouds to overcome appearance bias. For adapting large-scale models, Parameter-Efficient Fine-Tuning (PEFT) methods like LoRA \cite{hu2021lora} allow for rapid task-specific optimization without re-training the entire backbone. Furthermore, online adaptation frameworks like DAgger \cite{ross2011dagger} provide a mechanism for human-in-the-loop (HIL) correction to address failure cases, a paradigm we explore within the constraints of our training infrastructure.
